\documentclass[12pt,a4paper]{article}
\usepackage[utf8]{inputenc}
\usepackage[spanish]{babel}
\usepackage[T1]{fontenc}
\usepackage{geometry}
\geometry{top=2cm, bottom=2cm, left=2.5cm, right=2.5cm}
\usepackage{xcolor}
\usepackage{listings}
\usepackage{enumitem}
\usepackage{titlesec}
\usepackage{fancyhdr}
\usepackage{hyperref}

\definecolor{codegreen}{rgb}{0,0.6,0}
\definecolor{codegray}{rgb}{0.5,0.5,0.5}
\definecolor{codepurple}{rgb}{0.58,0,0.82}
\definecolor{backcolour}{rgb}{0.95,0.95,0.92}

\lstset{
    language=C,
    backgroundcolor=\color{backcolour},
    commentstyle=\color{codegreen},
    keywordstyle=\color{blue},
    numberstyle=\tiny\color{codegray},
    stringstyle=\color{codepurple},
    basicstyle=\ttfamily\small,
    breakatwhitespace=false,
    breaklines=true,
    captionpos=b,
    keepspaces=true,
    numbers=left,
    numbersep=5pt,
    showspaces=false,
    showstringspaces=false,
    showtabs=false,
    tabsize=2,
    frame=single
}

\pagestyle{fancy}
\fancyhf{}
\lhead{Sistemas Operativos 2026-I}
\rhead{Prueba Diagnóstica}
\cfoot{\thepage}

\title{\huge Prueba Diagnóstica: Introducción a la Programación en C}
\author{Monitorías de Sistemas Operativos}
\date{Marzo 2026}

\begin{document}

\maketitle

\section*{Instrucciones}
\begin{enumerate}
    \item Resuelva todos los ejercicios utilizando el lenguaje C.
    \item El código debe compilar correctamente.
    \item Organice su código de forma clara y utilice funciones cuando sea necesario.
    \item No use IA generativa (ChatGPT, Gemini, Copilot, etc.).
    \item Comente brevemente las partes importantes del programa si lo considera necesario.
\end{enumerate}
\newpage

\section*{Ejercicio 1 - Control de flujo}

Escriba un programa que lea dos números enteros $a$ y $b$. El programa debe recorrer todos los números enteros en el intervalo $[a, b]$ y para cada número $i$ debe imprimir:

\begin{itemize}
    \item Si $1 \le i \le 9$, imprima su representación en inglés en minúsculas (one, two, ..., nine).
    \item Si $i > 9$ y el número es par, imprima: \texttt{even}.
    \item Si $i > 9$ y el número es impar, imprima: \texttt{odd}.
\end{itemize}

\subsection*{Formato de entrada}
La entrada contiene dos líneas:
\begin{enumerate}
    \item $a$: Inicio del intervalo.
    \item $b$: Final del intervalo.
\end{enumerate}

\subsection*{Ejemplo}
\textbf{Entrada:}
\begin{lstlisting}[numbers=none]
8
11
\end{lstlisting}

\textbf{Salida:}
\begin{lstlisting}[numbers=none]
eight
nine
even
odd
\end{lstlisting}

\newpage

\section*{Ejercicio 2 - Vectores dinámicos y paso por referencia}

Escriba un programa que realice las siguientes tareas:
\begin{enumerate}
    \item Lea un número entero $n$ que representa la cantidad de elementos.
    \item Cree dinámicamente un vector de tamaño $n$ utilizando memoria dinámica (\texttt{malloc}).
    \item Lea $n$ números enteros desde la entrada estándar y guárdelos en el vector.
    \item Implemente una función que reciba el vector y calcule (usando punteros para devolver los valores):
    \begin{itemize}
        \item La suma de los elementos.
        \item El promedio.
        \item El valor máximo.
    \end{itemize}
\end{enumerate}

\subsection*{Requisitos}
\begin{itemize}
    \item El vector debe ser creado usando memoria dinámica.
    \item Debe utilizar punteros para devolver los resultados de la función (parámetros por referencia).
    \item Libere la memoria reservada al final del programa (\texttt{free}).
\end{itemize}

\subsection*{Ejemplo de ejecución}
\textbf{Entrada:}
\begin{lstlisting}[numbers=none]
5
10 4 8 2 6
\end{lstlisting}

\textbf{Salida:}
\begin{lstlisting}[numbers=none]
Sum: 30
Average: 6.00
Max: 10
\end{lstlisting}
\newpage

\section*{Ejercicio 3 - Lectura de archivos}

Se dispone de un archivo llamado \texttt{numbers.txt} que contiene un número entero por línea.

\subsection*{Tarea}
Escriba un programa en C que:
\begin{enumerate}
    \item Abra el archivo \texttt{numbers.txt} en modo lectura.
    \item Lea todos los números almacenados en el archivo.
    \item Almacene los valores en un vector dinámico.
    \item Calcule e imprima:
    \begin{itemize}
        \item La cantidad total de números.
        \item La suma de los números.
        \item El promedio.
    \end{itemize}
\end{enumerate}

\subsection*{Requisitos}
\begin{itemize}
    \item Usar \texttt{FILE*}, \texttt{fopen} y \texttt{fclose}.
    \item Usar \texttt{fscanf} o \texttt{fgets}.
    \item Utilizar memoria dinámica para el vector.
\end{itemize}


\end{document}
